\documentclass{article}
\usepackage[utf8x]{inputenc}
\usepackage[russian]{babel}
\usepackage{graphicx}
\title{Краткий отчёт о вычислениях производных функций}
\author{Postal Dude}
\date{November 2025}
\begin{document}
\maketitle
\newpage
\section{Вступление}
\hspace{5mm} Вступления пока нет, но оно скоро появится.
\section{Преобразования}
\begin{center}
Исходное выражение: \[ ( x ) + ( sin(  ( x ) * ( 2 ) ) ) \]\
\end{center}
Об этом мы сможем говорить в 11 классе, если живы будете

\begin{center}
\[\frac{d( 2 )}{dx}  =  0 \] \end{center}
Оценка на букву х - четыре

\begin{center}
\[\frac{d( x )}{dx}  =  1 \] \end{center}
В оригинале этой задачи этого условия не было, но поскольку вы слабы разумом, я его добавил

\begin{center}
\[\frac{d( ( x ) * ( 2 ) )}{dx}  =  ( ( 1 ) * ( 2 ) ) + ( ( x ) * ( 0 ) ) \] \end{center}
Оценка на букву х - четыре

\begin{center}
\[\frac{d( sin(  ( x ) * ( 2 ) ) )}{dx}  =  ( cos(  ( x ) * ( 2 ) ) ) * ( ( ( 1 ) * ( 2 ) ) + ( ( x ) * ( 0 ) ) ) \] \end{center}
Это ДИКИЙ интеграл, не надо его писать. Я его один раз напишу

\begin{center}
\[\frac{d( x )}{dx}  =  1 \] \end{center}
Оценка на букву х - четыре

\begin{center}
\[\frac{d( ( x ) + ( sin(  ( x ) * ( 2 ) ) ) )}{dx}  =  ( 1 ) + ( ( cos(  ( x ) * ( 2 ) ) ) * ( ( ( 1 ) * ( 2 ) ) + ( ( x ) * ( 0 ) ) ) ) \] \end{center}


\begin{figure}[h!]
\centering
\includegraphics[width=0.8\textwidth]{diffdump6.png}
\caption{График исходной функции}
\end{figure}



\begin{figure}[h!]
\centering
\includegraphics[width=0.8\textwidth]{diffdump7.png}
\caption{График производной функции}
\end{figure}


\section{Вывод}
Концовки пока нет, но она скоро появится
\\\\\\\\
Список литературы на лето и не только:\\\\
1) Л. Г. Петерсон "Математика. 2 класс. Углублённый уровень"\\
2) Лекция профессора А. С. Багирова "МАТЕМАТИКА - ЛЖЕНАУКА И ИНСТРУМЕНТ"ВЛИЯНИЯ ЯЩЕРСКОЙ ЦИВИЛИЗАЦИИ": \\ https://www.youtube.com/watch?v=V6G3sPbgubY \\
3) С. Лем "Солярис"\\
\end{document}